\chapter{Definición de la Organización y Requerimientos de Red}

\section{Descripción de la Organización}

El \textbf{Policlínico de Especialidades Médicas ``Salud \& Vida'' S.A.C.} es una institución de salud privada de nivel ambulatorio que brinda atención multidisciplinaria especializada. Se ubica en Tacna, Perú, y opera en una infraestructura de \textbf{planta única de 450~m$^2$ (30.00~m $\times$ 15.00~m)}.

\subsection{Misión y Perfil Tecnológico}

La organización requiere una red de datos de alta disponibilidad para soportar la gestión digital de historiales clínicos electrónicos (EMR/HCE), la transferencia de imágenes radiológicas en estándar DICOM, el control de inventarios de farmacia y la comunicación interna mediante telefonía IP (VoIP).

\begin{figure}[H]
\centering
\begin{tikzpicture}[
    scale=0.92, transform shape,
    box/.style={draw, rounded corners=3pt, fill=mdpiBlue!12, minimum width=2.8cm, minimum height=0.7cm, align=center, font=\footnotesize\bfseries, text=mdpiBlue!90!black},
    jefe/.style={draw, rounded corners=3pt, fill=mdpiBlue, minimum width=3.8cm, minimum height=0.75cm, align=center, font=\small\bfseries, text=white},
    linea/.style={-Stealth, thick, color=mdpiBlue!70}
]
\node[jefe] (dir) at (0,0) {Dirección General};
\node[box] (adm) at (-4.2,-1.6) {Administración\\Jefatura};
\node[box] (med) at (0,-1.6) {Área Médica\\(8 consultorios)};
\node[box] (diag) at (4.2,-1.6) {Diagnóstico\\Rayos~X / Lab.};
\node[box] (rec) at (-3.5,-3.3) {Recepción\\y Caja};
\node[box] (farm) at (0,-3.3) {Farmacia\\Asistencial};
\node[box] (triaje) at (3.5,-3.3) {Triaje\\y Tópico};
\draw[linea] (dir) -- (adm);
\draw[linea] (dir) -- (med);
\draw[linea] (dir) -- (diag);
\draw[linea] (adm) -- (rec);
\draw[linea] (med) -- (farm);
\draw[linea] (med) -- (triaje);
\end{tikzpicture}
\caption{Organigrama funcional del Policlínico ``Salud \& Vida'' S.A.C.}
\label{fig:organigrama}
\end{figure}

\subsection{Número de Trabajadores y Dispositivos}

La plantilla operativa asciende a \textbf{38 a 50 personas} entre médicos, enfermeras, técnicos, personal administrativo y personal de apoyo. El número de dispositivos de red que operan simultáneamente varía entre \textbf{44 y 50 terminales activos}, dentro del rango exigido por la rúbrica del proyecto (30 a 60 dispositivos).

\section{Inventario de Ambientes y Distribución de Planta}

La edificación se organiza en tres zonas funcionales claramente diferenciadas, tal como se muestra en el plano arquitectónico interactivo diseñado en SVG a escala 1:100:

\subsection{Zona Norte: Consultorios y Servicios Médicos}

\begin{table}[H]
\centering
\caption{Inventario de ambientes — Ala Norte}
\begin{tabularx}{\textwidth}{@{} >{\raggedright\arraybackslash}p{5.5cm} c c X @{}}
\toprule
\textbf{Ambiente} & \textbf{Dimensiones} & \textbf{Área} & \textbf{Función Principal} \\
\midrule
MDF (Telecomunicaciones) & 3.50~m $\times$ 4.20~m & 14.7~m$^2$ & Cuarto central de comunicaciones \\
Recepción / Admisión / Caja & 6.00~m $\times$ 4.20~m & 25.2~m$^2$ & Atención al público y facturación \\
Consultorio 01 & 3.50~m $\times$ 4.20~m & 14.7~m$^2$ & Medicina General \\
Consultorio 02 & 3.50~m $\times$ 4.20~m & 14.7~m$^2$ & Pediatría \\
Consultorio 03 & 3.50~m $\times$ 4.20~m & 14.7~m$^2$ & Ginecología \\
Consultorio 04 & 3.50~m $\times$ 4.20~m & 14.7~m$^2$ & Cardiología \\
Consultorio 05 & 3.50~m $\times$ 4.20~m & 14.7~m$^2$ & Traumatología \\
Farmacia Asistencial & 3.50~m $\times$ 4.20~m & 14.7~m$^2$ & Dispensación de medicamentos \\
\bottomrule
\end{tabularx}
\label{tab:ambientes-norte}
\end{table}

\subsection{Zona Central: Pasillo y Circulación}

El pasillo central de \textbf{30.00~m $\times$ 2.40~m (72~m$^2$)} constituye el eje troncal del sistema de canalización. Aloja la bandeja portacables tipo malla de 200$\times$50~mm suspendida a 2.75~m de altura, y los puntos de acceso Wi-Fi~6 AP-02 y AP-03 anclados al cielo raso.

\subsection{Zona Sur: Diagnóstico, Especialidades y Administración}

\begin{table}[H]
\centering
\caption{Inventario de ambientes — Ala Sur}
\begin{tabularx}{\textwidth}{@{} >{\raggedright\arraybackslash}p{5.5cm} c c X @{}}
\toprule
\textbf{Ambiente} & \textbf{Dimensiones} & \textbf{Área} & \textbf{Función Principal} \\
\midrule
Sala de Espera General & 7.00~m $\times$ 7.00~m & 49.0~m$^2$ & Sala pública para 35 personas \\
Triaje y Tópico & 3.50~m $\times$ 4.20~m & 14.7~m$^2$ & Registro y urgencias \\
Consultorio 06 & 3.50~m $\times$ 4.20~m & 14.7~m$^2$ & Oftalmología \\
Consultorio 07 & 3.50~m $\times$ 4.20~m & 14.7~m$^2$ & Odontología \\
Consultorio 08 & 3.50~m $\times$ 4.20~m & 14.7~m$^2$ & Psicología / Nutrición \\
Laboratorio Clínico & 4.50~m $\times$ 4.20~m & 18.9~m$^2$ & Bioquímica y análisis clínicos \\
Diagnóstico por Imágenes (Rayos~X) & 5.00~m $\times$ 4.20~m & 21.0~m$^2$ & Radiología y ecografía \\
Administración / Jefatura & 3.50~m $\times$ 7.00~m & 24.5~m$^2$ & Gestión administrativa y contable \\
\bottomrule
\end{tabularx}
\label{tab:ambientes-sur}
\end{table}

\section{Justificación de la Ubicación Geocéntrica del MDF}

El \textbf{Cuarto de Telecomunicaciones (MDF)} se sitúa en la posición central del edificio, con coordenadas georreferenciadas en el plano SVG de \(X = 14.00~\text{m},\; Y = 2.10~\text{m}\). Esta decisión reduce la distancia de la tirada horizontal más larga a \textbf{34.80~m} (Laboratorio Clínico), lo que representa apenas el \textbf{38.7\% del límite máximo de 90~m} establecido por la norma ANSI/TIA-568, otorgando un margen de seguridad superior al 61\%.

\section{Marco de Requerimientos de Red (Sección 6)}

\subsection{Cobertura de Red Inalámbrica (WLAN Wi-Fi 6)}

Se proyectan \textbf{cuatro Puntos de Acceso (APs) Wi-Fi~6 (IEEE~802.11ax)} PoE+ con radio efectivo de 8.0~m en banda de 5~GHz, distribuidos estratégicamente:

\begin{itemize}[leftmargin=*]
    \item \textbf{AP-01:} Sala de Espera General — red pública de cortesía para pacientes y pantalla de turnos.
    \item \textbf{AP-02:} Pasillo Central Oeste — cobertura médica a Consultorios 1 al 4 y Triaje.
    \item \textbf{AP-03:} Pasillo Central Este — cobertura médica a Consultorios 5 al 8 y Laboratorio.
    \item \textbf{AP-04:} Administración General (sector este) — cobertura corporativa a jefatura y contabilidad.
\end{itemize}

\subsection{Matriz de Requerimientos de Red por Ambiente}

\begin{table}[H]
\centering
\caption{Matriz de requerimientos de red por ambiente}
\small
\begin{tabularx}{\textwidth}{@{} c >{\raggedright\arraybackslash}p{4.2cm} X c >{\centering\arraybackslash}p{2.4cm} @{}}
\toprule
\textbf{N°} & \textbf{Ambiente} & \textbf{Dispositivos Asignados} & \textbf{Tomas} & \textbf{Medio Físico} \\
\midrule
1  & Admisión, Caja y Recepción & 3 PCs, 1 impresora, 2 VoIP & 6 & UTP Cat~6 \\
2  & Sala de Espera General    & 1 Smart TV, 1 AP-01 PoE & 2 & UTP Cat~6 \\
3  & Triaje y Tópico           & 2 PCs, 1 teléfono IP & 3 & UTP Cat~6 \\
4--8 & Consultorios 1 al 5     & 5 PCs, 5 teléfonos IP & 10 & UTP Cat~6 \\
9--11 & Consultorios 6 al 8   & 3 PCs, 3 teléfonos IP & 6 & UTP Cat~6 \\
12 & Laboratorio Clínico       & 2 PCs, 1 analizador bioquímico & 3 & F/UTP Cat~6A \\
13 & Rayos~X y Ecografía       & 2 estaciones DICOM, 1 impresora & 3 & F/UTP Cat~6A \\
14 & Farmacia Asistencial      & 2 terminales venta, 1 imp. térmica & 3 & UTP Cat~6 \\
15 & Administración y Jefatura & 4 PCs, 1 multifuncional, 1 VoIP, 1 AP & 7 & UTP Cat~6 \\
16 & Pasillo Central (APs)     & AP-02, AP-03 PoE Wi-Fi~6 & 2 & UTP Cat~6 \\
17 & MDF (Servidores)          & Servidor HCE, Servidor PACS, NAS & 3 & F/UTP Cat~6A \\
\midrule
   & \textbf{Subtotal Activos}  &  & \textbf{48} & \\
   & Reserva técnica ($>$15\%) &  & \textbf{8}  & \\
\midrule
   & \textbf{TOTAL PARCHEO}     &  & \textbf{56} & \\
\bottomrule
\end{tabularx}
\label{tab:requerimientos-red}
\end{table}

\subsection{Resumen de Características Técnicas Identificadas}

\begin{itemize}[leftmargin=*]
    \item \textbf{Rango de dispositivos simultáneos:} 44 a 50 terminales activos (dentro del rango de 30 a 60 exigido).
    \item \textbf{Tráfico crítico de imágenes médicas:} Archivos DICOM de 50~MB a 300~MB por estudio radiológico.
    \item \textbf{Tráfico VoIP:} Sensible a latencia y jitter; requiere calidad de servicio (QoS).
    \item \textbf{Zonas con EMI elevada:} Rayos~X (tubo de 100~kV) y Laboratorio (centrífugas y analizadores) exigen cable blindado F/UTP Cat~6A y puesta a tierra bajo ANSI/TIA-607-C.
    \item \textbf{Seguridad ante incendios:} Forro LSZH obligatorio en entornos hospitalarios (IEC~60332-1).
\end{itemize}
